% ================================================
% 文件: 02_常用测量仪器.tex
% 章节: 第18章 常用测量仪器
% 所属部分: 仪器仪表
% 版本: v1.0
% ================================================

\chapter{常用测量仪器}
\label{chap:common_instruments}

\section{示波器}
\label{sec:oscilloscope}

示波器是一种用于观察和测量电信号波形的仪器。与电压表、频率计等只能给出单一数值的仪器不同，示波器把随时间变化的电压以二维图形（横轴为时间、纵轴为电压）实时呈现出来，因此被称为“电子工程师的眼睛”。在通信设备的调试中，无论是检查时钟边沿、测量串口波形，还是观察调制包络，示波器都是首选工具。

\subsection{工作原理}

示波器通过高速电子束在荧光屏上扫描，将电信号转换为可见的波形。在数字示波器中，这一过程被分解为“采样—量化—存储—显示”四个步骤：模拟前端对被测信号进行衰减或放大，使其落入 ADC 的输入范围；ADC 按固定时钟对信号采样，把连续波形变成一串离散的数字样点；样点写入高速存储器（采集存储器）；显示引擎再以荧光屏或液晶屏重现波形。模拟示波器则依靠阴极射线管中受水平与垂直偏转板控制的电子束在磷光屏上描迹。

理解示波器性能，关键在于三个相互关联的技术指标：\textbf{带宽}、\textbf{采样率}与\textbf{存储深度}。带宽 $B$（通常指 $-3\,\mathrm{dB}$ 幅度点）决定了仪器能不失真通过的最高信号频率；采样率 $f_s$ 必须至少为被测信号最高频率分量的两倍（奈奎斯特准则）才能无混叠地重建波形，实际工程中常取 $4\sim10$ 倍以上；存储深度则决定了在给定采样率下能捕获多长的时间窗口，三者满足
\[
    \text{存储深度} = f_s \times \text{记录时长} .
\]
此外，示波器对阶跃响应的上升时间 $t_r$ 与带宽之间近似满足
\[
    t_r \approx \frac{0.35}{B} \quad (B\ \text{以 Hz 为单位，高斯响应近似}) ,
\]
例如 $100\,\mathrm{MHz}$ 带宽对应的上升时间约为 $3.5\,\mathrm{ns}$。测量快速数字边沿时，若示波器本身上升时间远大于被测信号，则测得值将显著偏慢。

\subsection{主要功能}

\begin{itemize}
    \item \textbf{波形显示}：观察信号的形状和变化。通过适当的垂直灵敏度（V/格）与时间基准（s/格）设置，可将毫伏级慢变化信号或高速数字脉冲都置于屏幕内观察。
    \item \textbf{电压测量}：测量信号的幅度和峰值。利用光标或自动测量可读出峰峰值 $V_{\mathrm{pp}}$、幅值、平均值、均方根值（RMS）等。真有效值对非正弦波尤为关键。
    \item \textbf{频率测量}：测量信号的频率和周期。现代数字示波器通过对一个或数个周期的时间戳计数，或由波形重建后求周期，给出频率与占空比。
    \item \textbf{相位测量}：测量两个信号之间的相位差。使用双通道示波器，量取两路波形过零点（或同相位点）间的时间差 $\Delta t$ 与周期 $T$，相位差
    \[
        \varphi = 360^\circ \times \frac{\Delta t}{T} .
    \]
    \item \textbf{脉冲测量}：测量脉冲的上升时间、下降时间、脉宽与过冲等参数。这些参数对数字电路与雷达脉冲的时序分析至关重要。
\end{itemize}

数字示波器还普遍支持\textbf{自动测量}、\textbf{多通道数学运算}（如两通道相减以观察差分信号）、\textbf{频谱显示}（对采集数据进行 FFT）以及\textbf{波形存储与回放}。利用 X--Y 显示模式，可将通道 1 接 X、通道 2 接 Y，绘制李萨如（Lissajous）图形，由图形切点数判断两信号的频率比与相位关系，在无频率计的条件下粗略测定未知频率。

\subsection{类型}

\begin{itemize}
    \item \textbf{模拟示波器}：基于阴极射线管工作。实时性强、无死区时间，但无法存储、分辨力受限于人眼判读，已逐步退出主流。
    \item \textbf{数字示波器}：将信号数字化后显示。具有存储、自动测量、接口通信等优势，是当今台式与手持示波器的主流。
    \item \textbf{存储示波器}：可以存储和回放波形。强调长记录与分段存储能力，便于捕捉偶发瞬态与单次事件。
    \item \textbf{数字荧光示波器}：具有高刷新率和波形累积功能。通过“数字荧光”技术以亮度梯度表现信号出现概率，便于发现偶发毛刺与抖动分布。
\end{itemize}

下表列出了选购与使用示波器时最需关注的几项指标及其物理含义。

\begin{table}[htbp]
    \centering
    \caption{示波器关键指标及其含义}
    \begin{tabular}{|p{3.2cm}|p{4.6cm}|p{6.2cm}|}
        \hline
        \textbf{指标} & \textbf{符号/典型值} & \textbf{含义与影响} \\
        \hline
        \textbf{带宽} & $B$（如 $100\,\mathrm{MHz}$） & 可不失真通过的最高频率；限制可测上升时间 \\
        \hline
        \textbf{采样率} & $f_s$（如 $1\,\mathrm{GS/s}$） & 每秒样点数；须 $\ge 2$ 倍信号最高频率 \\
        \hline
        \textbf{存储深度} & 如 $1\,\mathrm{M}$ 点 & 决定记录时长；深存储利于捕获长事件 \\
        \hline
        \textbf{垂直分辨力} & 8/10/12 bit & ADC 位数；决定幅度量化细度 \\
        \hline
        \textbf{输入阻抗} & $1\,\mathrm{M}\Omega\parallel\!15\,\mathrm{pF}$ 或 $50\,\Omega$ & 影响对被测电路的负载效应 \\
        \hline
    \end{tabular}
\end{table}

使用探头时务必进行\textbf{探头补偿}：将 $10\times$ 衰减探头接至示波器的校准方波输出端，调节探头上的补偿电容，使方波边沿既不过冲也不圆角，否则高频测量将出现系统性失真。

\section{信号发生器}
\label{sec:signal_generator}

信号发生器是用于产生各种波形电信号的仪器。它为被测设备或电路提供“已知的激励”，再通过示波器、频谱仪等观测“响应”，是构建一切测量实验的基础。在通信设备调试中，信号发生器常用于模拟天线端口的射频输入、产生本振信号或扫频激励。

\subsection{主要功能}

\begin{itemize}
    \item \textbf{波形产生}：产生正弦波、方波、三角波等波形。函数发生器提供基本波形，任意波形发生器（AWG）还能由用户编辑任意形状的波形并回放。
    \item \textbf{频率调节}：调节输出信号的频率。频率准确度与稳定度由内部时基（通常为石英晶振）决定，用日老化率、频率准确度（如 $\pm 1\times10^{-6}$）表征。
    \item \textbf{幅度调节}：调节输出信号的幅度。常以开路电压或匹配负载（$50\,\Omega$）上的电平表示，可用 $\mathrm{dBm}$ 或 $\mathrm{V_{pp}}$ 给出。
    \item \textbf{调制功能}：实现调幅、调频等调制方式。高级信号源还支持调相、脉冲调制、数字调制（如 QAM、PSK）及扫频输出，用于接收机灵敏度与选择性测试。
\end{itemize}

对于正弦输出，其瞬时值可写为
\[
    u(t) = U_m \sin(2\pi f t + \varphi_0) ,
\]
其中 $U_m$ 为幅度、$f$ 为频率、$\varphi_0$ 为初相。当信号源以电平 $P$（单位 $\mathrm{dBm}$）向 $50\,\Omega$ 负载输出时，对应电压有效值为
\[
    U_{\mathrm{rms}} = \sqrt{P\,\left[\mathrm{mW}\right] \times 50\,\Omega} ,
\]
例如 $0\,\mathrm{dBm}$（即 $1\,\mathrm{mW}$）对应 $U_{\mathrm{rms}} \approx 0.224\,\mathrm{V}$。注意信号源标称的电平常指 $50\,\Omega$ 匹配时的值，若负载失配则实际加到负载的功率会偏离标称值。

\subsection{类型}

\begin{itemize}
    \item \textbf{函数发生器}：产生标准函数波形。以 DDS（直接数字合成）技术为主流的现代函数发生器，频率分辨率高、切换快、相位连续。
    \item \textbf{任意波形发生器}：可以产生自定义波形。内部存储用户编辑的样点序列，经 DAC 重建输出，适用于模拟复杂真实信号。
    \item \textbf{信号源}：用于提供特定频率和幅度的信号。在射频领域常指“射频信号源”，输出频率可达微波波段，并具备精密的电平与调制控制。
    \item \textbf{脉冲发生器}：产生脉冲信号。用于数字电路时序测试、雷达与激光触发，强调上升时间、脉宽与重复频率的精确可调。
\end{itemize}

\section{经典学生信号源 J2465}
\label{sec:student_signal_source}

J2465型学生信号源是20世纪80年代中国广泛使用的教学仪器，由新金师范教学仪器厂生产。它是电子技术课程中不可或缺的实验设备，为一代电子爱好者和工科学生提供了入门学习的工具。在那个集成电路与数字仪器尚不普及的年代，J2465 以纯模拟电路实现多波形输出，结构直观、价格低廉、易于维护，因而长期出现在中学与中专的实验室课桌上。

\subsection{技术参数}

\begin{itemize}
    \item \textbf{频率范围}：10 Hz ~ 1 MHz（分多个频段）
    \item \textbf{输出波形}：正弦波、方波、三角波
    \item \textbf{输出幅度}：0 ~ 5 V（连续可调）
    \item \textbf{输出阻抗}：600 $\Omega$
    \item \textbf{电源}：220 V AC
\end{itemize}

上述指标属于早期教学仪器的典型水平：以 RC 振荡或 LC 振荡产生信号，经分档频段开关与微调旋钮改变频率，整体频率准确度远不及现代晶振基准，但足以满足滤波器、放大器、运放电路等基础实验的激励需求。

\subsection{面板结构}

J2465的控制面板包含以下主要部件：

\begin{itemize}
    \item \textbf{频率选择开关}：选择不同的频率范围。分档切换使各频段内频率覆盖合理，避免单一旋钮跨度过大。
    \item \textbf{频率微调旋钮}：在选定范围内精确调节频率。与频率选择开关配合，可在档内连续改变振荡回路参数。
    \item \textbf{波形选择开关}：选择正弦波、方波或三角波。通过切换输出网络或整形电路实现不同波形输出。
    \item \textbf{幅度调节旋钮}：调节输出信号的幅度。改变输出衰减或增益，使输出端得到 0~5 V 范围内所需电平。
    \item \textbf{输出端子}：提供信号输出接口。标称 600 $\Omega$ 输出阻抗，在使用时应尽量使负载与之匹配以减小波形失真。
    \item \textbf{电源开关}：控制仪器电源。接通后通常需短暂预热使振荡频率稳定。
\end{itemize}

\subsection{使用方法}

使用J2465学生信号源的基本步骤：

\begin{enumerate}
    \item 将仪器连接到220V交流电源
    \item 根据实验需求选择合适的波形
    \item 设置频率范围并微调至所需频率
    \item 调节输出幅度至合适大小
    \item 将输出端子连接到被测电路或示波器
\end{enumerate}

实际操作中建议先用示波器监视输出，确认波形形状与幅度符合预期后再接入被测电路；由于输出阻抗为 600 $\Omega$，测量时应留意负载对幅度的分压影响。

\subsection{历史意义}

J2465学生信号源在中国电子教育史上具有重要意义：

\begin{itemize}
    \item \textbf{普及教育}：为广大中学生和大学生提供了接触电子技术的机会。它将抽象的“频率”“波形”“幅度”概念变成可触摸、可观测的实验对象。
    \item \textbf{培养人才}：许多电子工程师和无线电爱好者都是通过它入门的。借助这台仪器，无数学习者第一次完成振荡、放大、整形等基础电路实验。
    \item \textbf{历史价值}：见证了中国电子工业从起步到发展的历程。作为国产教学仪器的代表，它承载了一代人的技术记忆，也折射出当时国内仪器制造业的水平与特点。
\end{itemize}

\section{万用表}
\label{sec:multimeter}

万用表是一种多功能测量仪器，可以测量电压、电流、电阻等多种电量。因其“一表多用”且价格相对低廉，万用表是电子与通信工程师随身必备的工具。现代数字万用表（DMM）在电压、电阻测量上已达相当高的精度，但使用者若不了解其输入特性与量程限制，仍可能得到误导性的读数。

\subsection{主要功能}

\begin{itemize}
    \item \textbf{电压测量}：测量直流和交流电压。交流电压档有“平均值响应”与“真有效值（True RMS）”之分，真有效值表对失真波形的读数才准确。
    \item \textbf{电流测量}：测量直流和交流电流。测量电流需将表笔串入回路，内部经分流电阻或电流传感器取样，故接入会改变原电路。
    \item \textbf{电阻测量}：测量电阻值。由内部恒流源供电测压而得，测在线电阻时须先断电并考虑并联支路影响。
    \item \textbf{二极管测试}：检测二极管的好坏。以恒流源加正向电压，正常硅管约显示 $0.5\sim0.7\,\mathrm{V}$ 的正向压降。
    \item \textbf{通断测试}：检测电路是否导通。当电阻低于设定阈值时蜂鸣报警，便于快速查线。
\end{itemize}

万用表测交流电压时，若采用真有效值转换，则对任意波形的有效值满足
\[
    U_{\mathrm{rms}} = \sqrt{ \frac{1}{T} \int_0^T u^2(t)\,dt } ;
\]
普通平均值表只对正弦波准确，遇到方波、含谐波波形会显著偏离。另一方面，万用表对被测电路存在\textbf{负载效应}：其直流电压档输入电阻通常为 $10\,\mathrm{M}\Omega$ 或更高，但早期或廉价表可能仅数百 $\mathrm{k}\Omega$。当测量高内阻信号源（如 $100\,\mathrm{k}\Omega$ 分压点）时，有限输入电阻会分压，引入加载误差，此时应选用高输入阻抗档位或改用示波器（探头 $10\times$ 时约 $10\,\mathrm{M}\Omega$）。

\subsection{类型}

\begin{itemize}
    \item \textbf{模拟万用表}：使用指针显示测量结果。由磁电式表头与分流/分压网络构成，无需电源即可测电流电压，但分辨力与防过载能力较弱。
    \item \textbf{数字万用表}：使用数字显示测量结果。以 ADC 为核心，分辨力高、读数直观、带自动极性，已成为主流。
    \item \textbf{钳形表}：可以测量交流电流而不需要断开电路。利用电流互感器感应导线电流，适合大电流场合与在线测量。
    \item \textbf{兆欧表}：用于测量高电阻。内部提供高压（如 $500\,\mathrm{V}$、 $1000\,\mathrm{V}$）以测量绝缘电阻，评估设备与安全接地间的绝缘水平。
\end{itemize}

使用万用表尤其要注意\textbf{安全类别（CAT 等级）}：CAT III、CAT IV 用于配电侧的测量，要求表笔与内部电路具备更高的瞬态过电压耐受与隔离能力。在通信机房与电力相关场合，务必使用符合相应 CAT 等级且经过认证的仪表，并正确选择量程与表笔，防止触电与电弧伤害。

\section{频率计}
\label{sec:frequency_counter}

频率计是用于测量信号频率的仪器。它以“在已知时间内数出周期数”这一直接而精确的方式工作，是标定振荡器、检验时钟与测量信号源频率准确度的标准工具。在通信领域，频率计常用来验证本振、晶振与发射载波的频率偏差。

\subsection{工作原理}

频率计通过计数一定时间内信号的周期数来计算频率。其核心是一条“闸门”：在闸门时间 $T_g$ 内，让被测信号经整形后驱动计数器累加脉冲数 $N$，则频率为
\[
    f = \frac{N}{T_g} .
\]
闸门时间越长，可累计的脉冲越多，分辨力越高——频率计的理论分辨力（量化误差）为
\[
    \Delta f = \frac{1}{T_g} ,
\]
例如 $T_g = 1\,\mathrm{s}$ 时分辨力为 $1\,\mathrm{Hz}$，$T_g = 10\,\mathrm{s}$ 时分辨力可达 $0.1\,\mathrm{Hz}$。这正是“以时间换精度”的典型体现。

\subsection{主要功能}

\begin{itemize}
    \item \textbf{频率测量}：测量信号的频率。是频率计最基本的功能，精度主要由内部时基的稳定度决定。
    \item \textbf{周期测量}：测量信号的周期。通过对一个周期内的高频标准时钟计数实现，适用于低频信号的精密测量。
    \item \textbf{计数功能}：对脉冲信号进行计数。在闸门时间内累计外部事件数，用于产量统计或事件计数。
    \item \textbf{时间间隔测量}：测量两个信号之间的时间间隔。由启动与停止两个通道分别触发，对标准时钟计数得到，常用于测量传播延时与脉冲宽度。
\end{itemize}

频率计的准确度由内部\textbf{晶振时基}的准确度与老化率决定：短期稳定度影响读数抖动，长期老化与温度漂移决定系统偏差。高精度频率计采用恒温晶振（OCXO）甚至铷钟作为时基，或外接更高等级的频率标准。下表概括了频率计的主要技术关注点。

将上述常用仪器组合起来，即可搭建一套基础的电子测量工位：以信号源提供激励、示波器观测波形、万用表查点电平、频率计锁定时钟。在连接时应特别注意\textbf{阻抗与端接}：射频与脉冲测量多用 $50\,\Omega$ 端接以防反射；高阻探头需做补偿校准；接地线过长会引入环路电感，使高频测量出现振铃。\textbf{安全}方面，测量市电或高压电路须使用相应 CAT 等级的表笔与探头，先选大量程再逐步缩小，切忌带电切换或超量程测量。随着总线化与计算机控制的普及，这些仪器还可通过 USB、LAN 或 GPIB 纳入自动测试系统，由软件统一调度、采集与判定，大幅提升测试效率与一致性。掌握它们的原理与恰当用法，是通信设备研发、生产与维护的基本功。

\begin{table}[htbp]
    \centering
    \caption{频率计主要技术指标}
    \begin{tabular}{|p{3.4cm}|p{4.8cm}|p{5.8cm}|}
        \hline
        \textbf{指标} & \textbf{典型范围/值} & \textbf{说明} \\
        \hline
        \textbf{测量范围} & 数 Hz ~ 数百 MHz 或更高 & 由输入放大与预分频决定 \\
        \hline
        \textbf{闸门时间} & $0.01\sim10\,\mathrm{s}$ 可选 & 越长分辨力越高 \\
        \hline
        \textbf{分辨力} & $1/T_g$ & 由闸门时间决定 \\
        \hline
        \textbf{时基准确度} & $\pm(0.1\sim10)\times10^{-6}$ & 决定频率测量系统误差 \\
        \hline
        \textbf{灵敏度} & 数十 mV$\sim$百 mV & 可测最小输入幅度 \\
        \hline
    \end{tabular}
\end{table}
